\documentclass{amsart}
% For dealing with references we use the comment environment
\usepackage{verbatim}
\newenvironment{reference}{\comment}{\endcomment}
%\newenvironment{reference}{}{}
\newenvironment{slogan}{\comment}{\endcomment}
\newenvironment{history}{\comment}{\endcomment}

% For commutative diagrams we use Xy-pic
\usepackage[all]{xy}

% We use 2cell for 2-commutative diagrams.
\xyoption{2cell}
\UseAllTwocells

% We use multicol for the list of chapters between chapters
\usepackage{multicol}

% This is generally recommended for better output
\usepackage{lmodern}
\usepackage[T1]{fontenc}

% For cross-file-references

% Package for hypertext links:
\usepackage{hyperref}

% For any local file, say "hello.tex" you want to link to please

% Theorem environments.
%
\theoremstyle{plain}
\newtheorem{theorem}[subsection]{Theorem}
\newtheorem{proposition}[subsection]{Proposition}
\newtheorem{lemma}[subsection]{Lemma}

\theoremstyle{definition}
\newtheorem{definition}[subsection]{Definition}
\newtheorem{example}[subsection]{Example}
\newtheorem{exercise}[subsection]{Exercise}
\newtheorem{situation}[subsection]{Situation}

\theoremstyle{remark}
\newtheorem{remark}[subsection]{Remark}
\newtheorem{remarks}[subsection]{Remarks}

\numberwithin{equation}{subsection}

% Macros
%
\def\lim{\mathop{\mathrm{lim}}\nolimits}
\def\colim{\mathop{\mathrm{colim}}\nolimits}
\def\Spec{\mathop{\mathrm{Spec}}}
\def\Hom{\mathop{\mathrm{Hom}}\nolimits}
\def\Ext{\mathop{\mathrm{Ext}}\nolimits}
\def\SheafHom{\mathop{\mathcal{H}\!\mathit{om}}\nolimits}
\def\SheafExt{\mathop{\mathcal{E}\!\mathit{xt}}\nolimits}
\def\Sch{\mathit{Sch}}
\def\Mor{\mathop{\mathrm{Mor}}\nolimits}
\def\Ob{\mathop{\mathrm{Ob}}\nolimits}
\def\Sh{\mathop{\mathit{Sh}}\nolimits}
\def\NL{\mathop{N\!L}\nolimits}
\def\CH{\mathop{\mathrm{CH}}\nolimits}
\def\proetale{{pro\text{-}\acute{e}tale}}
\def\etale{{\acute{e}tale}}
\def\QCoh{\mathit{QCoh}}
\def\Ker{\mathop{\mathrm{Ker}}}
\def\Im{\mathop{\mathrm{Im}}}
\def\Coker{\mathop{\mathrm{Coker}}}
\def\Coim{\mathop{\mathrm{Coim}}}

% Boxtimes
%
\DeclareMathSymbol{\boxtimes}{\mathbin}{AMSa}{"02}

%
% Macros for moduli stacks/spaces
%
\def\QCohstack{\mathcal{QC}\!\mathit{oh}}
\def\Cohstack{\mathcal{C}\!\mathit{oh}}
\def\Spacesstack{\mathcal{S}\!\mathit{paces}}
\def\Quotfunctor{\mathrm{Quot}}
\def\Hilbfunctor{\mathrm{Hilb}}
\def\Curvesstack{\mathcal{C}\!\mathit{urves}}
\def\Polarizedstack{\mathcal{P}\!\mathit{olarized}}
\def\Complexesstack{\mathcal{C}\!\mathit{omplexes}}
% \Pic is the operator that assigns to X its picard group, usage \Pic(X)
% \Picardstack_{X/B} denotes the Picard stack of X over B
% \Picardfunctor_{X/B} denotes the Picard functor of X over B
\def\Pic{\mathop{\mathrm{Pic}}\nolimits}
\def\Picardstack{\mathcal{P}\!\mathit{ic}}
\def\Picardfunctor{\mathrm{Pic}}
\def\Deformationcategory{\mathcal{D}\!\mathit{ef}}

\setlength{\emergencystretch}{3em}
\begin{document}
\title[Stacks Project, Tag 06XL]{317 --- Essential surjectivity of the derived functor to a Serre quotient}
\maketitle
\noindent Copyright (C) 2005--2025 Johan de Jong. Stacks Project authors. Source: \href{https://stacks.math.columbia.edu/tag/06XL}{Tag 06XL}.

\noindent Editorial difficulty note (not author text). Difficulty H2: The proof strictifies a complex in the quotient category in both directions: pushouts handle nonnegative degrees and pullbacks handle negative degrees. It must then replace terms by finite sums of images and intersections of kernels so the lifted differentials square to zero in the original category. This is an actual lifting construction for unbounded quotient complexes..

\noindent Editorial provenance note (not author text). History: excerpt from revision \texttt{a0444\allowbreak 6e57e\allowbreak c1fbc\allowbreak 252a8\allowbreak 71afc\allowbreak ec775\allowbreak 2fb28\allowbreak 07b14}, derived.tex, lines 6003--6100. Reference presentation was adapted; original-author context and core support blocks were retained or added. The mathematical wording is unchanged. Licensed under GNU FDL 1.2 or later; no invariant sections or cover texts. See ../licenses/COPYING. Context blocks reproduce author definitions and supporting results; they are not additional counted problems.

\begin{lemma}
\label{homology-lemma-serre-subcategory-is-kernel}
Let $\mathcal{A}$ be an abelian category.
Let $\mathcal{C} \subset \mathcal{A}$ be a Serre subcategory.
There exists an abelian category $\mathcal{A}/\mathcal{C}$
and an exact functor
$$
F : \mathcal{A} \longrightarrow \mathcal{A}/\mathcal{C}
$$
which is essentially surjective and whose kernel is $\mathcal{C}$.
The category $\mathcal{A}/\mathcal{C}$ and the functor $F$ are
characterized by the following universal property: For any exact
functor $G : \mathcal{A} \to \mathcal{B}$ such that
$\mathcal{C} \subset \Ker(G)$ there exists a factorization
$G = H \circ F$ for a unique exact functor
$H : \mathcal{A}/\mathcal{C} \to \mathcal{B}$.
\end{lemma}

\begin{lemma}
\label{lemma-derived-of-quotient}
Let $\mathcal{A}$ be an abelian category.
Let $\mathcal{B} \subset \mathcal{A}$ be a Serre subcategory.
Then $D(\mathcal{A}) \to D(\mathcal{A}/\mathcal{B})$
is essentially surjective.
\end{lemma}

\begin{proof}
We will use the description of the category $\mathcal{A}/\mathcal{B}$
in the proof of
Homology, Lemma \ref{homology-lemma-serre-subcategory-is-kernel}.
Let $(X^\bullet, d^\bullet)$ be a complex of $\mathcal{A}/\mathcal{B}$.
This means that $X^i$ is an object of $\mathcal{A}$ and
$d^i : X^i \to X^{i + 1}$ is a morphism in $\mathcal{A}/\mathcal{B}$
such that $d^i \circ d^{i - 1} = 0$ in $\mathcal{A}/\mathcal{B}$.

\medskip\noindent
For $i \geq 0$ we may write $d^i = (s^i, f^i)$ where $s^i : Y^i \to X^i$
is a morphism of $\mathcal{A}$ whose kernel and cokernel are in $\mathcal{B}$
(equivalently $s^i$ becomes an isomorphism in the quotient category)
and $f^i : Y^i \to X^{i + 1}$ is a morphism of $\mathcal{A}$.
By induction we will construct a commutative diagram
$$
\xymatrix{
& (X')^1 \ar@{..>}[r] & (X')^2 \ar@{..>}[r] & \ldots \\
X^0 \ar@{..>}[ru] &
X^1 \ar@{..>}[u] &
X^2 \ar@{..>}[u] &
\ldots \\
Y^0 \ar[u]_{s^0} \ar[ru]_{f^0} &
Y^1 \ar[u]_{s^1} \ar[ru]_{f^1} &
Y^2 \ar[u]_{s^2} \ar[ru]_{f^2} &
\ldots
}
$$
where the vertical arrows $X^i \to (X')^i$ become isomorphisms
in the quotient category. Namely, we first let
$(X')^1 = \Coker(Y^0 \to X^0 \oplus X^1)$ (or rather the
pushout of the diagram with arrows $s^0$ and $f^0$) which gives the
first commutative diagram. Next, we take
$(X')^2 = \Coker(Y^1 \to (X')^1 \oplus X^2)$. And so on.
Setting additionally $(X')^n = X^n$ for $n \leq 0$ we see that the map
$(X^\bullet, d^\bullet) \to ((X')^\bullet, (d')^\bullet)$
is an isomorphism of complexes in $\mathcal{A}/\mathcal{B}$.
Hence we may assume $d^n : X^n \to X^{n + 1}$ is given
by a map $X^n \to X^{n + 1}$ in $\mathcal{A}$ for $n \geq 0$.

\medskip\noindent
Dually, for $i < 0$ we may write $d^i = (g^i, t^{i + 1})$ where
$t^{i + 1} : X^{i + 1} \to Z^{i + 1}$ is an isomorphism in the
quotient category and $g^i : X^i \to Z^{i + 1}$ is a morphism.
By induction we will construct a commutative diagram
$$
\xymatrix{
\ldots &
Z^{-2} &
Z^{-1} &
Z^0 \\
\ldots &
X^{-2} \ar[u]_{t_{-2}} \ar[ru]_{g_{-2}} &
X^{-1} \ar[u]_{t_{-1}} \ar[ru]_{g_{-1}} &
X^0 \ar[u]_{t^0} \\
\ldots &
(X')^{-2} \ar@{..>}[u] \ar@{..>}[r] &
(X')^{-1} \ar@{..>}[u] \ar@{..>}[ru]
}
$$
where the vertical arrows $(X')^i \to X^i$ become isomorphisms
in the quotient category. Namely, we take
$(X')^{-1} = X^{-1} \times_{Z^0} X^0$. Then we take
$(X')^{-2} = X^{-2} \times_{Z^{-1}} (X')^{-1}$. And so on.
Setting additionally $(X')^n = X^n$ for $n \geq 0$ we see that the map
$((X')^\bullet, (d')^\bullet) \to (X^\bullet, d^\bullet)$
is an isomorphism of complexes in $\mathcal{A}/\mathcal{B}$.
Hence we may assume $d^n : X^n \to X^{n + 1}$ is given
by a map $d^n : X^n \to X^{n + 1}$ in $\mathcal{A}$
for all $n \in \mathbf{Z}$.

\medskip\noindent
In this case we know the compositions $d^n \circ d^{n - 1}$
are zero in $\mathcal{A}/\mathcal{B}$. If for $n > 0$ we replace
$X^n$ by
$$
(X')^n = X^n/\sum\nolimits_{0 < k \leq n} \Im(\Im(X^{k - 2} \to X^k) \to X^n)
$$
then the compositions $d^n \circ d^{n - 1}$ are zero for $n \geq 0$.
(Similarly to the second paragraph above we obtain an isomorphism of
complexes
$(X^\bullet, d^\bullet) \to ((X')^\bullet, (d')^\bullet)$.)
Finally, for $n < 0$ we replace $X^n$ by
$$
(X')^n = \bigcap\nolimits_{n \leq k < 0}
(X^n \to X^k)^{-1}\Ker(X^k \to X^{k + 2})
$$
and we argue in the same manner to get a complex in $\mathcal{A}$
whose image in $\mathcal{A}/\mathcal{B}$ is isomorphic to the given one.
\end{proof}
\end{document}
